\section{谱不变量的定义及其公理}
\label{sec:spectral-invariants-definition-axioms}

本节\emph{不}假设 $H$ 已归一化，但对任意给定哈密顿量，以 $\underline H$ 表示其归一化：
\begin{equation}
 \underline H(t,x)=H(t,x)-\frac1{\operatorname{vol}_\omega(M)}
 \int_MH(t,x)\,\mu_\omega.
 \label{eq:21-hamiltonian-normalization}
\end{equation}

\subsection{谱不变量的定义}

对每个给定的齐次量子上同调类 $a\in QH^*(M)$，记
$a^\flat=a_H^\flat\in HF_*(H,J)$ 为其在同构 $\Phi_H$ 下的像。记
\[
 i^\lambda:HF_*^\lambda(H,J)\longrightarrow HF_*(H,J)
\]
为由典范包含诱导的同态。

\begin{definition}
\label{def:21-spectral-invariant}
设 $H$ 是非退化的一周期哈密顿量，$J$ 是 $H$-正则的。对任意给定的
$0\neq a\in QH^*(M)$，考虑对子 $(H,J)$ 的弗洛尔闭链
$\alpha\in\ker\partial_{(H,J)}\subset CF_*(H)$，使其代表 $a^\flat$。定义
\[
 \rho((H,J);a):=\inf_{\alpha;\,[\alpha]=a^\flat}\lambda_H(\alpha),
\]
等价地，
\[
 \rho((H,J);a):=
 \inf\{\lambda\in\mathbb R\mid a^\flat\in\operatorname{Im}i^\lambda
 \subset HF_*(H,J)\}.
\]
\end{definition}

下面是关于谱数非平凡性的基本定理。厄舍 (Ush08) 在弗洛尔同调的抽象表述中给出
了 $\rho((H,J);a)$ 有限性的纯代数证明。这里的证明改编自 (Oh05c)、(Oh06a) 与
(Ush08) 中相应结论的证明。

\begin{theorem}
\label{thm:21-spectral-invariant-finite-independent}
设 $H$ 非退化，$0\neq a\in QH^*(M)$。
\begin{enumerate}[label=(\arabic*)]
 \item 对任意 $H$-正则的 $J$，$\rho((H,J);a)>-\infty$。
 \item $\rho((H,J);a)$ 的定义与 $H$-正则的 $J$ 的选取无关；公共值记为
 $\rho(H;a)$。
\end{enumerate}
\end{theorem}

\begin{proof}
由定理~\ref{thm:21-boundary-depth-properties}(2) 及边界深度的定义，
\begin{equation}
 \partial(CF_*(H,J))\cap CF_*^\lambda(H,J)
 \subset\partial(CF_*^{\lambda+\lVert H\rVert}(H,J)).
 \label{eq:21-boundary-depth-control}
\end{equation}
以下命题是弗洛尔理论极小极大法中的重要成分，对应于经典临界点理论中“极小极大
集合”的关键\emph{\termLinkingProperty{}}（这种观点例如见 (Bn82)）。

\begin{proposition}
\label{prop:21-floer-minimax-lower-bound}
设 $\alpha$ 是 $CF_*(H,J)$ 中的 $\partial$-闭链，且其同调类
$[\alpha]\neq0$，其中 $\partial=\partial_{(H,J)}$。则
\begin{equation}
 \inf_{\beta\in CF_*(H,J)}\{\lambda_H(\alpha-\partial\beta)\}>-\infty.
 \label{eq:21-minimax-finite}
\end{equation}
\end{proposition}

\begin{proof}
反设存在 $\beta_j\in CF_*(H,J)$ 使
$\lambda_H(\alpha-\partial\beta_j)\to-\infty$。记
\begin{equation}
 \lambda_j=\lambda_H(\alpha-\partial\beta_j).
 \label{eq:21-lambda-sequence}
\end{equation}
取子列后可设 $\lambda_j>\lambda_{j+1}$。归纳构造序列 $\beta'_k$，使
\begin{align}
 \lambda_H(\alpha-\partial\beta'_k)&\leq\lambda_k,
 \label{eq:21-beta-prime-cycle-bound}\\
 \lambda_H(\beta'_{k+1}-\beta'_k)&\leq\lambda_k+\lVert H\rVert.
 \label{eq:21-beta-prime-cauchy-bound}
\end{align}
一旦构成，令 $\alpha_N=\alpha-\partial\beta'_N$。由
\eqref{eq:21-beta-prime-cycle-bound} 与 $\lambda_j\to-\infty$，有
$[\alpha_N]=[\alpha]$ 且 $\alpha_N$ 在 $\lambda_H$ 非阿基米德拓扑中趋于 $0$。
另一方面，\eqref{eq:21-beta-prime-cauchy-bound} 表明 $\beta'_j$ 在
$CF_*(H,J)$ 的 $\lambda_H$ 非阿基米德拓扑中收敛。记极限为 $\beta'_\infty$。
因 $\partial$ 关于该拓扑连续，$\partial\beta'_N\to\partial\beta'_\infty$，故
$\alpha=\partial\beta'_\infty$，与 $[\alpha]\neq0$ 矛盾。

余下只须构造 $\beta'_j$。令 $\beta'_1=\beta_1$。我们将证明：在 $\beta'_k$ 已满足
\eqref{eq:21-beta-prime-cycle-bound} 的假设下，下述构造给出的 $\beta'_{k+1}$
满足把 $k$ 换成 $k+1$ 后的 \eqref{eq:21-beta-prime-cycle-bound}。由于
\[
 \lambda_H(\alpha-\partial\beta'_k)\leq\lambda_k,
 \quad \lambda_H(\alpha-\partial\beta_{k+1})=\lambda_{k+1}<\lambda_k,
\]
故 $\lambda_H(\partial(\beta_{k+1}-\beta'_k))\leq\lambda_k$。由
\eqref{eq:21-boundary-depth-control}，存在
$\gamma'_k\in CF_*^{\lambda_k+\lVert H\rVert}(H,J)$ 使
\[
 \partial(\beta_{k+1}-\beta'_k)=\partial\gamma'_k.
\]
定义 $\beta'_{k+1}=\beta'_k+\gamma'_k$，便有
$\lambda_H(\beta'_{k+1}-\beta'_k)=\lambda_H(\gamma'_k)\leq\lambda_k+\lVert H\rVert$。

\begin{lemma}
\label{lem:21-inductive-beta-prime}
链 $\beta'_{k+1}=\beta'_k+\gamma'_k$ 满足将 $k$ 换为 $k+1$ 后的
\eqref{eq:21-beta-prime-cycle-bound}，并满足
\eqref{eq:21-beta-prime-cauchy-bound}。
\end{lemma}

\begin{proof}
按 $\gamma'_k$ 的定义，
\begin{equation}
 \alpha-\partial\beta'_{k+1}
 =\alpha-\partial(\beta'_k+\gamma'_k)
 =\alpha-\partial\beta_{k+1}.
 \label{eq:21-beta-prime-induction}
\end{equation}
故 $\lambda_H(\alpha-\partial\beta'_{k+1})\leq\lambda_{k+1}$，即所需结论。
\end{proof}
命题得证。
\end{proof}

命题~\ref{prop:21-floer-minimax-lower-bound} 按定义立即推出任意
$0\neq a\in QH^*(M)$ 的 $\rho((H,J);a)$ 有限。下面证明它与 $J$ 无关。令 $J'$ 是
另一个 $H$-正则殆复结构，$j=\{J^s\}_{0\leq s\leq1}$ 满足 $J^0=J,J^1=J'$。
证明
\begin{equation}
 \rho((H,J');a)\leq\rho((H,J);a).
 \label{eq:21-spectral-j-inequality}
\end{equation}
交换 $J,J'$ 即得反向不等式。给定 $\epsilon>0$，选取 $(H,J)$ 的弗洛尔闭链
$\alpha$，使
\[
 \lambda_H(\alpha)<\rho((H,J);a)+\epsilon.
\]
考虑推论~\ref{cor:21-change-almost-complex-structure} 中的传递映射
$h_{(H,j)}:CF(H,J)\to CF(H,J')$。传递后的闭链
$\alpha'=h_{(H,j)}(\alpha)$ 满足
$\lambda_H(\alpha')\leq\lambda_H(\alpha)$。由定义又有
$\rho((H,J');a)\leq\lambda_H(\alpha')$，故
\[
 \rho((H,J');a)\leq\rho((H,J);a)+\epsilon.
\]
令 $\epsilon\to0$ 即得~\eqref{eq:21-spectral-j-inequality}，定理得证。
\end{proof}

\begin{remark}
\label{rem:21-abstract-finiteness-proof}
检查命题~\ref{prop:21-floer-minimax-lower-bound} 的证明可见，其中唯一使用具体弗洛尔
链复形 $CF_*(H,J)$ 的地方，是定理~\ref{thm:21-boundary-depth-properties}(2) 给出的
\stfhIndex{I042}{Bianjieshendu}{边界深度（boundary depth）}
边界深度有限性。事实上，一旦一般微分 $\partial$ 的边界深度有限，逐字照搬上述
证明即可在弗洛尔理论的抽象框架中证明定理 13.8.4 所述谱不变量的有限性。
\end{remark}

考虑任意接近线性同伦 $s\mapsto(1-s)H+sK$ 的连接 $H,K$ 的同伦，可证明：

\begin{proposition}
\label{prop:21-spectral-continuity}
对任意非退化 $H,K$，
\begin{equation}
 -E^+(H-K)\leq\rho(H;a)-\rho(K;a)\leq E^-(H-K).
 \label{eq:21-spectral-continuity-bound}
\end{equation}
特别地，$\rho_a:H\mapsto\rho(H;a)$ 在 $C^0$ 拓扑（或
$L^{(1,\infty)}$ 拓扑）中连续，因而可连续延拓到
$C^0([0,1]\times M;\mathbb R)$。
\end{proposition}

\begin{proof}
任取 $\delta>0$，选取 $H$ 的闭链 $\alpha$，使 $[\alpha]=a^\flat$ 且
\begin{equation}
 \lambda_H(\alpha)\leq\rho(H;a)+\delta.
 \label{eq:21-near-minimizing-cycle}
\end{equation}
这里使用了\emph{已经证明}的 $\rho(H;a)>-\infty$。取从 $H$ 到 $K$ 的线性同伦
$h_{HK}^{\mathrm{lin}}$ 的通有逼近，并用推论~\ref{cor:21-filtered-continuation-map}，
可找到连接 $H,K$ 的 $\mathcal H$，使
\[
 \lambda_K(h_\mathcal H(\alpha))
 \leq\lambda_H(\alpha)+\int_0^1-\min_x(K_t-H_t)\,dt+\delta.
\]
结合~\eqref{eq:21-near-minimizing-cycle}，右端不超过
\[
 \rho(H;a)+2\delta+\int_0^1-\min_x(K_t-H_t)\,dt.
\]
因 $[h_{HK}^{\mathrm{lin}}(\alpha)]=a^\flat$，定义给出
$\lambda_K(h_{HK}^{\mathrm{lin}}(\alpha))\geq\rho(K;a)$，所以
\[
 \rho(K;a)-\rho(H;a)
 \leq2\delta+\int_0^1-\min_x(K_t-H_t)\,dt.
\]
令 $\delta\to0$，再交换 $H,K$，得到
\[
 \rho(H;a)-\rho(K;a)
 \leq\int_0^1-\min_x(H_t-K_t)\,dt
 =\int_0^1\max_x(K_t-H_t)\,dt.
\]
于是
\[
 \int_0^1-\max_x(K_t-H_t)\,dt
 \leq\rho(K;a)-\rho(H;a)
 \leq\int_0^1-\min_x(K_t-H_t)\,dt,
\]
即~\eqref{eq:21-spectral-continuity-bound}。它显然使 $\rho$ 可由连续性延拓到任意
$L^{1,\infty}$ 哈密顿量。
\end{proof}

\subsection{谱不变量的公理}

以下以公理列表给出函数 $\rho$ 的基本性质。

\begin{theorem}
\label{thm:21-spectral-axioms}
设 $(M,\omega)$ 是任意闭辛流形。对任意 $0\neq a\in QH^*(M)$，存在连续函数
\footnote{译者注：定理先把 $\rho$ 的定义域写成整个 $QH^*(M)$，但前句及本节定义
只对 $0\neq a$ 定义 $\rho(H;a)$；第 (7) 项在 $a+b=0$ 时也会落入同一问题。
此外，第 (5) 项使用了未在定理量词中声明的 $b$。译文照录。}
\[
 \rho:C^\infty(S^1\times M;\mathbb R)\times QH^*(M)\longrightarrow\mathbb R,
\]
满足下列公理。设 $H,F\in\mathcal H_m$ 是光滑含时哈密顿函数，且
$a\neq0\in QH^*(M)$。
\begin{enumerate}[label=(\arabic*)]
 \item 若 $r:[0,1]\to\mathbb R$ 光滑，则
 $\rho(H+r;a)=\rho(H;a)-\int_0^1r(t)\,dt$。
 \item （射影不变性）对任意 $0\neq\lambda\in\mathbb Q$，
 $\rho(H;\lambda a)=\rho(H;a)$。
 \item （归一化）若 $a=\sum_{A\in\Gamma}a_Aq^{-A}$，则
 $\rho(0;a)=v(a)$，其中 $0$ 是零函数，且
 \[
  v(a):=\min\{\omega(-A)\mid a_A\neq0\}
  =-\max\{\omega(A)\mid a_A\neq0\}
 \]
 是 $a$ 的向上赋值。
 \item （辛不变性）对任意辛微分同胚 $\eta$，
 $\rho(\eta_*H;\eta_*a)=\rho(H;a)$。特别地，若
 $\eta\in\operatorname{Symp}_0(M,\omega)$，则 $\rho(\eta_*H;a)=\rho(H;a)$。
 \item （三角不等式）$\rho(H\#F;a\cdot b)\leq\rho(H;a)+\rho(F;b)$。
 \item （哈密顿连续性）
 $|\rho(H;a)-\rho(F;a)|\leq\lVert H\#\overline F\rVert
 =\lVert H-F\rVert$，其中 $\lVert\cdot\rVert$ 是
 $C^\infty(S^1\times M;\mathbb R)$ 上的 $L^{(1,\infty)}$-范数。特别地，
 $\rho_a:H\mapsto\rho(H;a)$ 是 $C^0$-连续的。
 \item （加法三角不等式）
 $\rho(H;a+b)\leq\max\{\rho(H;a),\rho(H;b)\}$。
\end{enumerate}
\end{theorem}

\begin{proof}
对第 (1) 项，$H+r$ 与 $H$ 有完全相同的哈密顿流、弗洛尔方程和相应离壳设定；
唯一差别是作用量泛函平移了 $-\int_0^1r(t)dt$，故结论来自谱不变量的构造。
射影不变性也由构造显然成立。$C^0$-连续性立即来自命题
\ref{prop:21-spectral-continuity}。三角不等式的证明推迟至 21.4 节。

证明辛不变性时，对任意辛微分同胚 $\eta:(M,\omega)\to(M,\omega)$ 考察辛共轭
\[
 \phi\longmapsto\eta\phi\eta^{-1};\qquad
 \operatorname{Ham}(M,\omega)\longrightarrow\operatorname{Ham}(M,\omega).
\]
回忆推前函数 $\eta_*H(t,x)=H(t,\eta^{-1}(x))$；若 $H\mapsto\phi$，它生成共轭
$\eta\phi\eta^{-1}$。与滤弗尔同调相关的这一共轭的基本事实如下：
\begin{enumerate}[label=(\arabic*)]
 \item 若 $H\mapsto\phi$，则 $\eta_*H\mapsto\eta\phi\eta^{-1}$；
 \item $H$ 非退化蕴含 $\eta_*H$ 非退化；
 \item $(H,J)$ 弗洛尔正则蕴含 $(\eta_*H,\eta_*J)$ 弗洛尔正则；
 \footnote{译者注：原文把推前后的\personFloer{}对子排成 $(\eta_*J,\eta_*H)$，与全文
“哈密顿量在前、殆复结构在后”的顺序相反；译文按类型次序改正。}
 \item 存在自然双射
 $\eta_*:\widetilde{\mathcal L}_0(M)\to\widetilde{\mathcal L}_0(M)$，
 $\eta_*([z,w])=[\eta\circ z,\eta\circ w]$，且
 $\mathcal A_H([z,w])=\mathcal A_{\eta_*H}(\eta_*[z,w])$；
 \item $\operatorname{Symp}(M,\omega)$ 自然作用于 $QH^*(M)$；若
 $\alpha=\sum_{[z,w]}a_{[z,w]}[z,w]$ 代表 $a^\flat$，则
 $\eta_*(\alpha)=\sum_{[z,w]}a_{[z,w]}\eta_*([z,w])$ 代表 $(\eta_*a)^\flat$；
 \item 相应作用量泛函的 $L^2$-梯度满足
 $\eta_*(\operatorname{grad}_J\mathcal A_H)([z,w])
 =\operatorname{grad}_{\eta_*J}(\mathcal A_{\eta_*H})(\eta_*[z,w])$；
 \item 若 $u:\mathbb R\times S^1\to M$ 是 $(H,J)$ 的扰动柯西--黎曼方程之解，
 则 $\eta_*u=\eta\circ u$ 是 $(\eta_*J,\eta_*H)$ 的解；两组数据及其解具有相同的
 一切弗雷德霍姆性质；
 \item $\eta\in\operatorname{Symp}(M,\omega)$ 以拉回 $a\mapsto\eta^*a$ 作用于
 $QH^*(M)$。\footnote{译者注：公理与第 (5) 项以 $\eta_*a$ 记量子上同调上的作用，第 (8) 项却
称其为拉回 $\eta^*a$；原文记号不一致，译文逐处保留。}
\end{enumerate}
这些事实说明，共轭诱导典范的保滤过链同构
\[
 \eta_*:(CF_*^\lambda(H),\partial_{(H,J)})\longrightarrow
 (CF_*^\lambda(\eta_*H),\partial_{(\eta_*H,\eta_*J)}),
\]
其中 $\lambda\in\mathbb R\setminus\operatorname{Spec}(H)
=\mathbb R\setminus\operatorname{Spec}(\eta_*H)$。因此它在同调上诱导保滤过同构
\[
 \eta_*:HF_*^\lambda(H,J)\longrightarrow HF_*^\lambda(\eta_*H,\eta_*J).
\]
辛不变性随即来自构造 $\rho(H;a)$ 所用的弗洛尔极小极大过程。
\end{proof}

以上构造 $\rho(H;a)$ 一直假设 $H$ 一周期；这是哈密顿弗洛尔同调构造的附带要求。
通过重参数化哈密顿道路这一简单的柔性技巧可以去除此要求。这样的延拓使谱不变量
在应用中非常灵活，并在连续哈密顿动力系统的各类应用中十分重要（例如见 (Oh10)）。
下面说明如何去掉周期性，并把 $\rho(H;a)$ 延拓到任意含时哈密顿量
$H:[0,1]\times M\to\mathbb R$。显然霍费尔范数 $\lVert H\rVert$ 的定义不要求周期性。

由\emph{未必周期}的 $H$ 构造 $H^\zeta$ 型含时周期哈密顿量，其中 $\zeta$ 是 $[0,1]$ 的
重参数化，满足
\begin{equation}
 \zeta(t)=
 \begin{cases}
 0,&0\leq t\leq\epsilon_0/2,\\
 1,&1-\epsilon_0/2\leq t\leq1,
 \end{cases}
 \label{eq:21-reparameterization-cutoff}
\end{equation}
以及 $\zeta'(t)\geq0$。定义
\[
 H^\zeta(t,x)=\zeta'(t)H(\zeta(t),x),
\]
它生成哈密顿同痕 $t\mapsto\phi_H^{\zeta(t)}$。任意两个这样的重参数化函数
$\zeta_0,\zeta_1$ 所给道路 $\phi_{H^{\zeta_0}}$、$\phi_{H^{\zeta_1}}$ 彼此同伦，故
$\rho(H^{\zeta_0};a)=\rho(H^{\zeta_1};a)$。这使我们能给出如下定义。

\begin{definition}
\label{def:21-nonperiodic-spectral-invariant}
设 $H$ 是任意含时、未必一周期的哈密顿量。定义
\begin{equation}
 \rho(H;a):=\rho(H^\zeta;a),
 \label{eq:21-nonperiodic-spectral-definition}
\end{equation}
其中 $\zeta$ 是任一这样的重参数化。
\end{definition}

利用归一化 $H$ 与相应哈密顿道路 $\phi_H:t\mapsto\phi_H^t$ 的一一对应，在道路空间
\[
 P^{\mathrm{ham}}(\operatorname{Symp}(M,\omega),\mathrm{id})
\]
上定义谱函数
\[
 \rho_a:P^{\mathrm{ham}}(\operatorname{Symp}(M,\omega),\mathrm{id})
 \longrightarrow\mathbb R.
\]

\begin{definition}[哈密顿道路的谱不变量]
\label{def:21-hamiltonian-path-spectral-invariant}
设 $\lambda$ 是哈密顿道路。定义
\[
 \rho(\lambda;a):=\rho(\operatorname{Dev}(\lambda);a),
\]
即若 $\lambda=\phi_H$，则 $\rho(\lambda;a)=\rho(\underline H;a)$。
\end{definition}

以 $\widetilde{\operatorname{Ham}}(M,\omega)$ 表示下述道路空间中道路相对于端点的
同伦类之集：
\[
 P^{\mathrm{ham}}(\operatorname{Symp}(M,\omega),\mathrm{id}).
\]
它就是 $\operatorname{Ham}(M,\omega)$ 的万有覆盖空间，并赋予商拓扑。一个重要
问题是：当 $\lambda,\mu$ 相对于端点道路同伦时，是否
\[
 \rho(\lambda;a)=\rho(\mu;a),
\]
从而 $\rho$ 能下降到 $\widetilde{\operatorname{Ham}}(M,\omega)$？稍后将在 21.5 节讨论。
\footnote{译者注：原文指向 21.5 节，但同伦不变性的正式陈述与证明实际位于
21.6 节；21.5 节提供证明所需的谱性。译文保留原交叉引用并登记。}
